\begin{longtable}{@{}P{1.0cm}P{4.6cm}P{3.6cm}P{0.3cm}P{0.3cm}P{0.3cm}P{0.3cm}P{0.3cm}P{0.3cm}P{0.3cm}P{0.3cm}P{0.3cm}P{0.36cm}P{0.75cm}P{3.0cm}P{3.9cm}@{}}
\caption{Data extraction (part 2): limitations, reproducibility, quality items Q1--Q10 (1 = yes, .5 = partial, 0 = no, -- = N/A; items defined in Section 7) and decision evidence.}\label{tab:ext2}\\
\toprule Study & Limitations (authors; ours) & Reproducibility & 1 & 2 & 3 & 4 & 5 & 6 & 7 & 8 & 9 & 10 & QA\% & QA notes & Decision evidence \\ \midrule \endfirsthead
\toprule Study & Limitations & Reproducibility & 1 & 2 & 3 & 4 & 5 & 6 & 7 & 8 & 9 & 10 & QA\% & QA notes & Decision evidence \\ \midrule \endhead
\bottomrule \endfoot
\textbf{S01} & Authors: robustness may conflict with other goals; enforcing robustness proposed as future work (Sec.4). Ours: short workshop paper. & Uses author implementations of LIME and SHAP and DeepExplain (p.3); no paper code link. & 1 & 1 & .5 & 1 & 1 & .5 & .5 & .5 & -- & .5 & 72 & Q3 dataset details brief; Q6 small-scale; Q8 libraries only; Q10 sampling described, no repeated seeds. & Abstract, Secs.2-3 (pp.1-4): local Lipschitz robustness of LIME and SHAP on UCI tabular data and COMPAS (plus images); read: arXiv:1806.08049 (ICML 2018 Workshop on Human Interpretability in Machine Learning) \\
\textbf{S02} & Ours: short workshop paper; descriptive evidence; no stability index. & No code link found. & 1 & 1 & .5 & .5 & .5 & .5 & .5 & 0 & -- & 1 & 61 & Q3 data preprocessing brief; Q4 no baseline; Q5 selection probability only; Q6 descriptive; Q8 none; Q10 100 repeated runs. & Abstract, Secs.3-4 (pp.1-4): three sources of uncertainty in LIME on synthetic tree data, COMPAS (tabular) and 20 Newsgroups (text); read: arXiv:1904.12991v2 (ICML 2019 AI for Social Good workshop paper) \\
\textbf{S03} & Ours: 5-page workshop paper; adversarial robustness rather than benign stability; no repetitions reported. & Code: github.com/\allowbreak{}seansaito/\allowbreak{}Faster-LIME (p.1). & 1 & .5 & .5 & 1 & .5 & .5 & .5 & 1 & -- & .5 & 67 & Q2/\allowbreak{}Q3 brief; Q5 detection accuracy only; Q6 small experiments; Q10 no repetitions. & Abstract, Secs.2-6 (pp.1-4): robustness of LIME to the adversarial scaffolding attack using CTGAN sampling on tabular COMPAS and two other datasets; read: arXiv:2006.12302v3 version of the AdvML '20 (KDD workshop) paper \\
\textbf{S04} & Authors: other explainer classes left for future work (p.7). Ours: adversarial manipulation, not benign perturbation stability. & Code: github.com/\allowbreak{}dylan-slack/\allowbreak{}Fooling-LIME-SHAP (p.1). & 1 & 1 & 1 & 1 & .5 & 1 & .5 & 1 & 1 & .5 & 85 & Q5 rank-share outcome; Q7 brief; Q9 discrimination and fairwashing discussed; Q10 repetitions not reported. & Abstract, Secs.3-5 (pp.1-6): scaffolding attack hides biased classifiers from LIME and SHAP on tabular COMPAS, Communities and Crime and German Credit; read: arXiv:1911.02508 version of the AIES 2020 paper (doi 10.1145/\allowbreak{}3375627.3375830) \\
\textbf{S05} & Authors: relies on local linear approximation; still vulnerable to sampling-based attacks if LIME/\allowbreak{}SHAP sampling is kept (Conclusion). & No code link found in the text read. & 1 & 1 & 1 & 1 & 1 & 1 & 1 & .5 & -- & 1 & 94 & Q8 no code link in text; Q10 repeated sampling with significance tests. & Abstract, Secs.3-5 (pp.1-9): BayesLIME/\allowbreak{}BayesSHAP credible intervals and local-Lipschitz stability on tabular COMPAS and German Credit (plus images); read: arXiv:2008.05030 version of the NeurIPS 2021 paper \\
\textbf{S06} & Authors: global importances only; one dataset confidential. Ours: local explanations not assessed; mostly regression datasets; statistical tests not reported. & UCI datasets (except proprietary); no code link found in the full text. & 1 & 1 & .5 & 1 & 1 & 1 & .5 & .5 & -- & 1 & 83 & Q3 proprietary dataset undisclosed; Q7 limitations brief; Q8 no code link; Q10 50 seeded repetitions. & Abstract, Sec.3 (pp.1,4-6): stability of TreeSHAP/\allowbreak{}gain importances of tree ensembles under input and model perturbation; read: arXiv:2110.00086 (pre-published version of the CD-MAKE 2021 paper, doi 10.1007/\allowbreak{}978-3-030-84060-0\_\allowbreak{}19) \\
\textbf{S07} & Ours: stability defined only as run-to-run feature-set variability; no input perturbation. & Code: github.com/\allowbreak{}rehmanzafar/\allowbreak{}dlime\_\allowbreak{}experiments. & 1 & 1 & 1 & 1 & 1 & 1 & .5 & 1 & -- & 1 & 94 & Q7 brief; Q10 repeated runs. & Abstract and experiments (HTML full text): deterministic LIME (DLIME) vs LIME stability on 6 public and 6 synthetic tabular datasets; read: Publisher HTML (Mach. Learn. Knowl. Extr. 3(3):525-541, 2021) \\
\textbf{S08} & Authors: controls only sampling variance, not sensitivity to kernel parameters or model variability; temporal data remain open (Sec.6). Ours: single-instance illustrations, small data, no input-perturbation analysis. & Code: github.com/\allowbreak{}ZhengzeZhou/\allowbreak{}slime (p.5); public datasets. & 1 & 1 & 1 & 1 & 1 & .5 & .5 & 1 & -- & 1 & 89 & Q6 few instances; Q7 scope limits noted; Q8 code and data public; Q9 no personal data beyond public benchmark (N/\allowbreak{}A); Q10 repetitions control sampling noise. & Abstract, Secs.4-5 (pp.1,4-6): LIME stability method evaluated on tabular breast-cancer and simulated data; read: arXiv:2106.07875 version of the KDD'21 paper (doi 10.1145/\allowbreak{}3447548.3467274) \\
\textbf{S09} & Authors: none stated beyond motivating limitations of earlier metrics. Ours: short workshop paper, ANN-focused representation metric, no tree ensembles, no statistical testing. & No code link in the paper; public datasets. & 1 & 1 & .5 & 1 & 1 & .5 & 0 & .5 & -- & .5 & 67 & Q3 datasets named, sizes not given; Q6 single figure of medians; Q7 no limitations section; Q8 data public, code not linked; Q10 number of perturbations/\allowbreak{}seeds not detailed. & Secs.3-4 (pp.2-5): relative stability of 7 explainers on tabular Adult, COMPAS, German Credit; read: arXiv:2203.06877 (preprint; presented at an ICLR 2022 workshop) \\
\textbf{S10} & Authors: indices guarantee concordance, not fidelity; framework relies on a ridge local model (Sec.7). Ours: single proprietary dataset, illustrative instances, no input perturbation. & Code: github.com/\allowbreak{}giorgiovisani/\allowbreak{}LIME\_\allowbreak{}stability (p.2); data proprietary. & 1 & 1 & .5 & .5 & 1 & .5 & .5 & .5 & .5 & 1 & 70 & Q3 dataset size not reported; Q4 no competing stability measure; Q6 case-study evidence; Q8 code yes, data no; Q9 regulation discussed; Q10 repeated calls isolate sampling noise. & Abstract, Secs.4-6 (pp.1,6-15): LIME stability indices on tabular credit-risk data; read: Preprint arXiv:2001.11757 (compiled 13 Nov 2020) of the JORS article \\
\textbf{S11} & Authors: building local models without model access is hard for image data (p.9). Ours: finite-sample uncertainty rather than input perturbation; preprint. & Code: github.com/\allowbreak{}surinahn/\allowbreak{}xai-uncertainty-quantification (p.1). & 1 & 1 & 1 & 1 & 1 & .5 & .5 & 1 & -- & 1 & 89 & Q6 real-data results visual only; Q7 brief; Q10 simulation over parameter grid. & Abstract, Secs.2-5 (pp.1-9): bootstrap uncertainty intervals for local explanations; simulation plus tabular COMPAS, German Credit and Spambase; read: Preprint arXiv:2301.05761v3 (24 Jun 2023, 'under review'; no peer-reviewed version found) \\
\textbf{S12} & Authors: limitations discussed in Suppl. O (p.11). & No code link found in the full text; public datasets. & 1 & 1 & 1 & 1 & 1 & 1 & .5 & .5 & -- & 1 & 89 & Q7 limitations in supplement; Q8 no code link; Q10 multiple environments and reported SDs. & Secs.4-5 (pp.5-10): LINEX stability metrics (CI, CAC) on tabular datasets (e.g., IRIS, MEPS) plus images; read: arXiv:2201.12143 version of the NeurIPS 2023 paper \\
\textbf{S13} & Authors: random forest only; dataset-dependence (p.10). Ours: perturbations not prediction-preserving. & Public datasets; no code link. & 1 & 1 & 1 & .5 & 1 & 1 & 1 & .5 & -- & .5 & 83 & Q4 one model; Q8 data only; Q10 seeds not reported. & Abstract, Secs.3-7 (pp.1-10): SHAP stability under shuffling, missingness and Gaussian noise at four levels on four tabular datasets; read: Zenodo preprint of the DSAA 2023 paper (doi 10.1109/\allowbreak{}DSAA60987.2023.10302535) \\
\textbf{S14} & Authors: single dataset and model family (Tiny Paper scope). Ours: neural model only (not tree ensembles); no input perturbation. & Code publicly accessible (footnote, p.1); MIMIC-III requires credentialed access. & 1 & 1 & 1 & .5 & 1 & 1 & .5 & 1 & .5 & 1 & 85 & Q4 compares background sizes, no alternative explainer; Q7 brief; Q8 code public; Q9 clinical data, ethics not discussed; Q10 100 resamplings per size. & Abstract, Secs.2-3 (pp.1-5): SHAP stability vs background-data size on tabular MIMIC-III; read: arXiv:2204.11351 (ICLR 2023 Tiny Paper version) \\
\textbf{S15} & Ours: short proceedings paper; stability metric briefly defined; no statistical tests. & No code link found. & 1 & .5 & 1 & .5 & .5 & .5 & .5 & 0 & -- & .5 & 56 & Q2 brief; Q4 LIME only; Q5 brief; Q8 none; Q10 not detailed. & Abstract, Secs.3-4 (pp.1-7): CTGAN-LIME stability vs LIME on Heart Disease, Red Wine Quality and Census Income; read: Publisher PDF (Proc. 10th World Congress EECSS'24, paper CIST 158, open access) \\
\textbf{S16} & Review scope: adversarial manipulation rather than benign perturbation; non-systematic search. & Paper list on GitHub (hbaniecki/\allowbreak{}adversarial-explainable-ai). & 1 & 1 & -- & -- & -- & 1 & 1 & 1 & 1 & -- & 100 & Review: dataset/\allowbreak{}baseline/\allowbreak{}metric items N/\allowbreak{}A; ethics (fairwashing) discussed. & Whole paper: survey of adversarial attacks/\allowbreak{}defences on explanations incl. SHAP/\allowbreak{}LIME on tabular models (review); read: arXiv:2306.06123 version of the Information Fusion survey (doi 10.1016/\allowbreak{}j.inffus.2024.102303) \\
\textbf{S17} & Authors: extension to images and text proposed (p.9). Ours: determinism trivially yields perfect run-to-run stability. & No code link found. & 1 & 1 & .5 & 1 & 1 & .5 & .5 & 0 & -- & 1 & 72 & Q3 dataset sizes brief; Q6 small; Q8 none; Q10 10 repeated explanations. & Abstract, Secs.3-4 (pp.1100-1108): ST-LIME fixed-noise sampling and Feature Sequence Stability Index vs five LIME variants on four tabular datasets; read: Publisher PDF (Int. J. Intelligent Engineering and Systems 17(6), 2024, open access) \\
\textbf{S18} & Authors: the correlated case needs accurate pre-computed matrices; LIME and distillation trees are also unstable (p.14). & Code: github.com/\allowbreak{}jeremy-goldwasser/\allowbreak{}ControlSHAP (p.2). & 1 & 1 & 1 & 1 & 1 & 1 & .5 & 1 & -- & 1 & 94 & Q7 limited; Q10 repeated estimation runs. & Abstract and Sec.5 (pp.1,9-10): Monte-Carlo variability of Shapley estimates reduced with control variates on tabular data; read: arXiv:2310.07672 version of the xAI 2024 paper (doi 10.1007/\allowbreak{}978-3-031-63797-1\_\allowbreak{}21) \\
\textbf{S19} & Authors: depends on decision-boundary estimation, trading computation time against accuracy (p.9). & Code: github.com/\allowbreak{}davinhill/\allowbreak{}GPEC (p.6). & 1 & 1 & 1 & 1 & .5 & .5 & 1 & 1 & -- & .5 & 83 & Q5 uncertainty estimate rather than a stability metric; Q6 largely visual comparisons; Q10 seeds not reported in sections read. & Abstract, Secs.3-5 (pp.1-9): GPEC uncertainty for KernelSHAP/\allowbreak{}BayesSHAP/\allowbreak{}BayesLIME explanations on tabular Census, Online Shoppers and German Credit (plus images); read: arXiv:2210.02419 version of the AISTATS 2024 paper (PMLR 238) \\
\textbf{S20} & Authors: single dataset; only size and imbalance examined (p.5). & Code: github.com/\allowbreak{}anushaihalapathirana/\allowbreak{}data\_\allowbreak{}characteristics\_\allowbreak{}shap; data restricted. & 1 & 1 & 1 & .5 & .5 & 1 & 1 & .5 & .5 & 1 & 80 & Q4 three models but no stability baseline; Q5 variance only; Q8 code, no data; Q9 health data, ethics not discussed; Q10 cross-validation folds. & Abstract and Secs.2-3 (pp.1-5): SHAP stability vs data size, imbalance and background size on tabular T2D EHR data; read: Publisher PDF (ESANN 2024 proceedings, pp.297-302) \\
\textbf{S21} & Authors: future work on other explainers. Ours: run-to-run stability only (no input perturbation). & Code: github.com/\allowbreak{}gwladyskelodjou/\allowbreak{}st-shap (p.4). & 1 & 1 & 1 & 1 & 1 & 1 & .5 & 1 & -- & 1 & 94 & Q7 limited; Q10 repeated executions. & Abstract and experiments (pp.1-8): Kernel SHAP run-to-run instability and a fully stable variant on tabular datasets; read: arXiv:2312.12115 version of the AAAI 2024 paper (doi 10.1609/\allowbreak{}aaai.v38i12.29208) \\
\textbf{S22} & Ours: two medical datasets; one model; variation from resampling, not input perturbation. & Not checked in sections read. & 1 & 1 & .5 & 1 & 1 & .5 & .5 & .5 & .5 & 1 & 75 & Q3 dataset sizes not checked; Q6 two case studies; Q8 not verified; Q9 health data; Q10 multiple k values. & Abstract, Secs.III-V (pp.1-2,4-5,13-14): generalisability, concordance and stability of SHAP and LIME combined with k-fold on tabular hypothyroidism and diabetes data; read: Publisher PDF (IEEE Access 2024, open access) \\
\textbf{S23} & Authors: bound could be tightened; tabular LIME with discretisation is not astute by this definition; robustness can conflict with correctness (p.9). & Public explainer implementations and OpenXAI used (pp.7-8); no paper-specific code link found. & 1 & 1 & 1 & 1 & 1 & 1 & 1 & .5 & -- & 1 & 94 & Q8 libraries public, no paper code; Q10 20 explainer runs. & Abstract, Secs.3-5 (pp.1-9): explainer astuteness (probabilistic Lipschitzness) for SHAP, RISE and CXPlain on simulated and UCI tabular data (plus MNIST); read: arXiv:2206.12481 version of the AISTATS 2024 paper (PMLR 238) \\
\textbf{S24} & Authors: future work on user studies and image/\allowbreak{}text data (p.15). Ours: no direct stability comparison with LIME or SHAP in the experiments read; robustness tested by resampling, not input perturbation. & Code: github.com/\allowbreak{}Moffran/\allowbreak{}calibrated\_\allowbreak{}explanations; pip/\allowbreak{}conda package (p.7). & 1 & 1 & 1 & .5 & .5 & 1 & .5 & 1 & -- & 1 & 83 & Q4 no competing explainer in stability tests; Q5 variance only; Q7 brief; Q10 30 runs per dataset. & Abstract, Secs.2 and 5 (pp.1,3,12-14): stability (30 seeded runs) and robustness (30 resampled refits) of Calibrated Explanations on 25 tabular datasets; read: arXiv:2305.02305 version (Semantic Scholar lists Expert Systems with Applications) \\
\textbf{S25} & Authors: robustness values in arbitrary units, comparison-only. Ours: preprint; tiny datasets; no statistical tests. & Public UCI datasets; no code link found. & 1 & 1 & 1 & 1 & .5 & .5 & .5 & .5 & -- & .5 & 72 & Q5 Lipschitz units hard to compare; Q6 small-scale; Q8 no code; Q10 sampling of neighbourhoods described but no repeated seeds. & Secs.2-3 (pp.2-6): SHAP robustness (local Lipschitz) for RF with prediction-preserving neighbourhoods on 4 UCI tabular datasets; read: Preprint arXiv:2402.19025 (not peer reviewed) \\
\textbf{S26} & Ours: single small dataset; preprint; perturbation design loosely specified. & Code: github.com/\allowbreak{}josesousaribeiro/\allowbreak{}XAI-eXirt-vs-Trust (p.9). & 1 & .5 & .5 & 1 & 1 & .5 & .5 & 1 & -- & .5 & 72 & Q2/\allowbreak{}Q3 limited detail; Q6 descriptive; Q7 none explicit; Q10 perturbation levels but no repetitions reported. & Abstract and Sec.5.2 (pp.1,7-13): stability of SHAP and other XAI rankings under test-data perturbation on tabular diabetes data; read: Preprint arXiv:2407.03108 (not peer reviewed) \\
\textbf{S27} & Authors: dependence on the Lipschitz constant; efficacy of stable rank needs more study (pp.2,17). Ours: neural networks only; toy tabular data. & No code link found. & 1 & 1 & .5 & 1 & 1 & .5 & 1 & 0 & -- & .5 & 72 & Q3 toy datasets; Q6 limited empirical scope; Q8 none; Q10 sampling not detailed. & Abstract and Sec.6 (pp.1,11-17): probabilistic-Lipschitz robustness of LIME/\allowbreak{}IG/\allowbreak{}SmoothGrad/\allowbreak{}SHAP with tabular XOR and Iris (plus MNIST); read: Preprint arXiv:2402.18863 (not peer reviewed) \\
\textbf{S28} & Ours: one dataset family; stability of the sampling only, no prediction-preserving input perturbation. & Code: github.com/\allowbreak{}zokaityte/\allowbreak{}LIME-Explainability-Stability (Data Availability). & 1 & 1 & 1 & 1 & 1 & .5 & .5 & 1 & -- & 1 & 89 & Q6 one dataset; Q7 brief; Q10 repeated LIME runs. & Abstract and results (HTML full text): LIME stability and fidelity with Weibull, Gamma, Beta and Pareto perturbation distributions on tabular CIC-IDS-2018; read: Publisher HTML (Mach. Learn. Knowl. Extr. 7(1):21, 2025) \\
\textbf{S29} & Authors: future work on cost distributions, synthetic data, alternative cost matrices and more stable XAI (Sec.6). Ours: instability is induced by retraining on resampled data, so explainer sampling noise (LIME) is not separated from model variation; no input-perturbation or prediction-preservation analysis. & Public datasets; supplementary material cited; no code link found in the full text. & 1 & 1 & 1 & 1 & 1 & 1 & .5 & .5 & .5 & .5 & 80 & Q7 future-work only; Q8 data public, code not linked; Q9 regulatory context (GDPR/\allowbreak{}AI Act) but no privacy analysis; Q10 25 iterations per rate but seeds and LIME randomness not separated. & Abstract \& Sec.3 (pp.1,4-8): SHAP/\allowbreak{}LIME stability on 4 tabular credit datasets is the stated aim; read: Accepted manuscript arXiv:2509.01409 (states published in EJOR, doi 10.1016/\allowbreak{}j.ejor.2025.05.039) \\
\textbf{S30} & Authors: online variance estimate uses one mini-batch; higher-order interactions not covered (p.10). & Code: github.com/\allowbreak{}nliulab/\allowbreak{}SIM-Shapley. & 1 & 1 & .5 & 1 & 1 & 1 & 1 & 1 & -- & 1 & 94 & Q3 datasets named, sizes in appendix; Q10 repeated runs measure estimator variance. & Abstract, Secs.2-4: stability (run-to-run consistency) of a Shapley estimator on tabular datasets; read: Preprint arXiv:2505.08198 ('under review') \\
\textbf{S31} & Authors: SPRT is conservative; future less conservative designs (p.9). Ours: neural models only; no input perturbation. & Code: github.com/\allowbreak{}jeremy-goldwasser/\allowbreak{}feature-rankings (p.2). & 1 & 1 & .5 & 1 & 1 & 1 & .5 & 1 & -- & 1 & 89 & Q3 dataset details in appendix; Q7 brief; Q10 repeated replicates with error-rate control. & Abstract, Sec.6 (pp.1,7-8): sampling instability of SHAP/\allowbreak{}LIME top-K rankings with hypothesis-testing fixes; tabular Adult Census among datasets; read: Preprint arXiv:2401.15800v5 (3 Jul 2025); DBLP lists an earlier version titled 'Provably Stable Feature Rankings with SHAP and LIME' \\
\textbf{S32} & Ours: stability operationalisation briefly described; single dataset; no repetitions or tests reported in sections read. & No code link; public dataset. & 1 & .5 & 1 & 1 & .5 & .5 & .5 & .5 & .5 & .5 & 65 & Q2/\allowbreak{}Q5 stability procedure brief; Q6 descriptive; Q9 legal defensibility discussed; Q10 not reported. & Abstract and Secs.3-4 (HTML full text): SHAP vs LIME on XGBoost and TabNet with explanation stability and Jaccard consistency metrics on tabular UNSW-NB15 (borderline); read: Publisher HTML (Applied Sciences 15(13):7329, 2025) \\
\textbf{S33} & Authors: relevant works may be missed given more than 20,000 citing papers (p.5); LIME only, no comparison with SHAP (p.2). Ours: selection process not reported as counts. & Online overview of LIME variants (Fig.1, p.4). & 1 & .5 & -- & -- & -- & 1 & 1 & 1 & -- & -- & 90 & Review: dataset/\allowbreak{}baseline/\allowbreak{}metric/\allowbreak{}randomness items N/\allowbreak{}A; Q2 selection criteria described but no counts. & Abstract, Secs.1-5 (pp.1-17): survey and taxonomy of LIME variants organised by fidelity, stability and domain challenges (review); read: arXiv:2503.24365 (Semantic Scholar lists the xAI 2025 conference version) \\
\textbf{S34} & Authors: single dataset, single model class, no hyperparameter tuning (p.13). Ours: global SHAP only; seeds are the only perturbation. & Public dataset; no code link found. & 1 & 1 & 1 & .5 & 1 & 1 & 1 & .5 & .5 & 1 & 85 & Q4 no comparison explainer or model; Q7 explicit limitations; Q8 data public, code not linked; Q9 regulatory (ECOA/\allowbreak{}FCRA) framing; Q10 100 seeds. & Secs.2-3 (pp.2-12): SHAP ranking consistency across 100 seed-varied XGBoost models on Taiwan credit data; read: Published article (Risks 2025) via arXiv:2508.01851 PDF \\
\textbf{S35} & Ours: SHAP/\allowbreak{}LIME only as comparators; stability under bootstrap resampling, not input perturbation. & No code link found in the sections read. & 1 & 1 & .5 & 1 & 1 & 1 & .5 & 0 & -- & 1 & 78 & Q3 dataset sizes not checked; Q7 method limitations only; Q8 none found; Q10 50 bootstrap resamples. & Abstract, Secs.I-IV (pp.1-3,7-8): bootstrap stability of Bayesian-AIME vs SHAP and LIME on tabular Titanic, breast-cancer and synthetic data; read: Publisher PDF (IEEE Access 2025, open access) \\
\textbf{S36} & Authors: undersampled reference, older CPRD release, no external validation (p.13). & CPRD data restricted; no code link found. & 1 & 1 & 1 & 1 & 1 & 1 & 1 & .5 & 1 & 1 & 95 & Q8 restricted data, no code; Q9 anonymised data, licence noted; Q10 100 bootstraps. & Methods 2.1-2.4, Results (pp.2-13): LIME/\allowbreak{}SHAP consistency under class imbalance on tabular CPRD EHR data; read: Publisher PDF (Frontiers in Artificial Intelligence 8:1682919, 2025) \\
\textbf{S37} & Ours: single proprietary dataset; perturbation design not prediction-preserving. & Data on request; no code link. & 1 & 1 & 1 & 1 & .5 & .5 & .5 & 0 & -- & .5 & 67 & Q5 consistency metric brief; Q6 one dataset; Q8 none; Q10 repetitions unclear. & Abstract and Secs.3-4 (HTML full text): Shapley/\allowbreak{}Least Core hybrid attribution improves explanation consistency under data perturbations for XGBoost rent prediction (tabular); read: Publisher HTML (Buildings 15(17):3133, 2025) \\
\textbf{S38} & Ours: single region and small sample; stability only across CV folds (retraining), not input perturbation. & Data contained in the article; no code link. & 1 & 1 & 1 & 1 & .5 & .5 & .5 & .5 & -- & 1 & 78 & Q5 CV only; Q6 single study area; Q7 brief; Q8 no code; Q10 10 folds. & Abstract and Conclusions (HTML full text): SHAP interpretive robustness of a hybrid LightGBM/\allowbreak{}XGBoost/\allowbreak{}RF model measured by coefficient of variation over 10-fold CV on tabular landslide factors; read: Publisher HTML (ISPRS Int. J. Geo-Inf. 14(7):277, 2025) \\
\textbf{S39} & Ours: simulated quantum hardware; run-to-run stability only. & Uses public lime and lime\_\allowbreak{}stability code and Qiskit; datasets public. & 1 & 1 & 1 & .5 & 1 & 1 & .5 & 1 & -- & 1 & 89 & Q4 no competing explainer; Q7 brief; Q10 repeated LIME calls. & Abstract and Secs.3-5 (HTML full text): LIME stability (VSI/\allowbreak{}CSI over repeated calls) for quantum SVMs and quantum neural networks on tabular UCI datasets; read: Publisher HTML (Machine Learning 115(6):141, 2026) \\
\textbf{S40} & Authors: US label measures disruption duration, not injury (abstract). Ours: perturbations not filtered to preserve predictions. & Code, split indices and artifacts: github.com/\allowbreak{}malhawarat/\allowbreak{}sae-xcrash. & 1 & 1 & 1 & 1 & 1 & 1 & 1 & 1 & -- & 1 & 100 & Q7 negative results reported; Q8 full release; Q10 audit criteria pre-specified. & Abstract and SHAP audit Step 4 (HTML full text): perturbation stability of SHAP (10 Gaussian-noise copies per instance) on tabular US-Accidents and UK STATS19 (borderline: one of four audit steps); read: Publisher HTML (Applied Sciences 16(12):5876, 2026) \\
\textbf{S41} & Authors: single in-distribution benchmark; transfer to other datasets untested (p.20). Ours: stability across resampling only; no input perturbation. & No code link found; public benchmark. & 1 & 1 & 1 & 1 & 1 & 1 & 1 & .5 & -- & 1 & 94 & Q8 data public, no code; Q10 ten seeds. & Secs.2.4, 3.4 (pp.8,16): TreeSHAP rank stability under training-pool resampling for XGBoost/\allowbreak{}RF on tabular CICIoT2023; read: Publisher PDF (Journal of Expert Systems and Sustainable Development 2(1):247-268, 2026; received/\allowbreak{}revised/\allowbreak{}accepted dates printed) \\
\textbf{S42} & Ours: data available on request only for processed files; single dataset suite. & Data available on request; no code link. & 1 & 1 & 1 & 1 & 1 & 1 & .5 & .5 & .5 & 1 & 85 & Q7 brief; Q8 on request; Q9 financial data, synthetic suite; Q10 repeated permutations and retraining. & Abstract and Methods (HTML full text): SHAP ranking stability across five biased/\allowbreak{}shifted Bank Account Fraud variants vs permutation importance; read: Publisher HTML (FinTech 5(3):77, 2026) \\
\textbf{S43} & Authors: the result concerns the single-model paradigm; conditional SHAP does not escape it (p.7). Ours: preprint; very long theoretical document. & Code and Lean proofs: github.com/\allowbreak{}DrakeCaraker/\allowbreak{}dash-impossibility-lean and dash-shap. & 1 & 1 & 1 & 1 & 1 & 1 & 1 & 1 & -- & 1 & 100 & Q10 seeds/\allowbreak{}trials and CIs reported. & Executive summary, Secs.3-5, 9 (pp.5-18,35): impossibility of stable SHAP rankings under collinearity; experiments on 11 datasets and a 77-dataset survey; read: Preprint arXiv:2605.21492 (not peer reviewed) \\
\textbf{S44} & Ours: stability mainly theoretical; group-level (not feature-level) attributions. & Data described in article; no code link. & 1 & 1 & 1 & 1 & .5 & 1 & .5 & .5 & -- & 1 & 83 & Q5 theoretical bound; Q7 brief; Q8 no code; Q10 20 seeds. & Title, abstract and Sec.3.1.9 (HTML full text): Lipschitz attribution-stability guarantee for TreeSHAP-importance-weighted group attributions; 7 tabular regression datasets (borderline); read: Publisher HTML (Mathematics 14(13):2419, 2026) \\
\textbf{S45} & Authors: several design choices (uniform prior, hard hypothesis matching) not claimed optimal (p.14). & Preprocessing and hyperparameters deferred to the final version; no code link. & 1 & .5 & .5 & .5 & 1 & .5 & 1 & 0 & -- & 1 & 67 & Q2/\allowbreak{}Q3 details deferred; Q4 limited baselines; Q6 small experiments; Q8 none. & Abstract, Secs.4-5 (pp.1,10-14): variability of LIME/\allowbreak{}SHAP via Weight of Evidence on tabular Titanic and synthetic data; read: Preprint arXiv:2609.00090 (not peer reviewed) \\
\textbf{S46} & Authors: controlled tabular classification only; design choices of the metric (p.11). & Code to be released upon acceptance (p.6). & 1 & 1 & 1 & 1 & 1 & 1 & 1 & 0 & -- & 1 & 89 & Q8 code not yet released; Q10 50 runs per instance. & Secs.3-5 (pp.3-11): run-to-run stability of LIME/\allowbreak{}SHAP rankings on four tabular datasets; read: Preprint arXiv:2608.02697 (accepted at EXPLAINS 2026 per arXiv listing) \\
\textbf{S47} & Authors: random-forest focus; reducing epistemic uncertainty needs better data (Summary). Ours: uncertainty rather than perturbation stability. & Code: github.com/\allowbreak{}dubeyakshat07/\allowbreak{}UBiQTree; Zenodo doi 10.5281/\allowbreak{}zenodo.17542114. & 1 & 1 & 1 & .5 & 1 & 1 & .5 & 1 & .5 & 1 & 85 & Q4 no competing UQ baseline reported in sections read; Q9 public health data, ethics brief; Q10 sampling-based. & Summary and Methods (pp.1,9-10): decomposes variability (uncertainty) of SHAP values for tree ensembles on tabular clinical datasets; read: arXiv:2508.09639 version of the Patterns article (doi 10.1016/\allowbreak{}j.patter.2025.101454) \\
\textbf{S48} & Authors: no local meteorological predictors; RPI does not establish causality (p.22). & Supplementary material online; no code link found. & 1 & 1 & 1 & 1 & 1 & 1 & 1 & .5 & -- & 1 & 94 & Q8 supplement only; Q10 bootstrap replicates. & Abstract, Secs.2-3 (pp.1-11): SHAP rank stability under station-level bootstrap, predictor perturbation and model class in a reliability performance index (tabular lagged-predictor data); read: Publisher PDF (Machine Learning: Science and Technology 7(4):045033, open access) \\
\textbf{S49} & Authors: limited to shallow networks (p.4). & No code link found; public datasets. & 1 & .5 & 1 & .5 & 1 & .5 & .5 & .5 & -- & 1 & 72 & Q2 brief (6-page format); Q4 no external baseline; Q6 small experiments; Q8 no code; Q10 repeated training runs. & Abstract and Sec.4 (pp.1,4-5): stability of Shapley-value explanations across a Rashomon set of near-optimal models on tabular data; read: Publisher PDF (ESANN 2026 proceedings, pp.145-150) \\
\textbf{S50} & Authors: single dataset; XGBoost quantised at explanation layer due to tooling (pp.5,11). Ours: no repeated runs or significance tests. & Data public (Kaggle); no code link. & 1 & 1 & 1 & 1 & 1 & 1 & 1 & .5 & -- & .5 & 89 & Q7 explicit limitations; Q8 data only; Q10 single trained model per family, no seeds reported. & Abstract, Results (pp.1,3-11): SHAP stability under post-training quantisation of LR/\allowbreak{}RF/\allowbreak{}XGBoost on tabular IEEE-CIS fraud data; read: Publisher PDF (Cureus Journal of Computer Science, published 4 Jun 2026; review dates printed) \\
\textbf{S51} & Ours: only three retraining seeds; no input perturbation. & Public datasets; no code link found. & 1 & 1 & 1 & 1 & 1 & 1 & .5 & .5 & -- & .5 & 83 & Q7 brief; Q8 data only; Q10 three seeds only. & Abstract and Sec.4 (HTML full text): retraining stability of SHAP rankings (seeds 42-44) within an XAI reliability audit on tabular UNSW-NB15 and IoT-ToN (borderline); read: Publisher HTML (J. Computer Virology and Hacking Techniques 22(1):80, 2026, open access) \\
\textbf{S52} & Authors: weak separation when uncertainty has low dispersion; no native uncertainty for LightGBM/\allowbreak{}CatBoost (pp.10,19). & Code: github.com/\allowbreak{}comrados/\allowbreak{}uq-xai (p.1). & 1 & 1 & 1 & 1 & 1 & 1 & 1 & 1 & -- & 1 & 100 & Q8 code public; Q10 perturbation types and budgets defined. & Abstract, Sec.4 (pp.1,6-9): explanation stability of four XAI methods on four UCI tabular datasets vs epistemic uncertainty; read: Preprint arXiv:2603.29915 (not peer reviewed) \\
\textbf{S53} & Ours: retraining stability only; regional health dataset. & Python packages listed; no code link. & 1 & 1 & 1 & .5 & 1 & 1 & .5 & .5 & .5 & 1 & 80 & Q4 two models only; Q8 no code; Q9 anonymised health data; Q10 repeated fits with fixed seeds. & Abstract, Methods (pp.113-120): SHAP ranking stability (Sequential Rank Agreement) of XGBoost/\allowbreak{}LightGBM across 64 simulated conditions and BPJS claims data; read: Publisher PDF via journal viewer (Jurnal Statistika dan Komputasi 5(1):113-, 2026, open access) \\
\textbf{S54} & Ours: significance argued from non-overlapping SD-based intervals; education domain only. & Not checked in sections read. & 1 & 1 & 1 & 1 & 1 & .5 & .5 & .5 & .5 & 1 & 80 & Q6 inference from SD intervals; Q8 not verified; Q9 student data, ethics not seen; Q10 repeated runs. & Abstract, Secs.III-IV (pp.1-2,5-11): LIME vs H-LIME stability (Jaccard, Spearman over 10 perturbation runs) on two tabular educational datasets; read: Publisher PDF (IEEE Access 2026, open access) \\
\textbf{S55} & Ours: single comparison of two scenarios; stability score loosely defined; no repetitions or tests. & No code or data link found. & 1 & .5 & .5 & 0 & .5 & 0 & .5 & 0 & .5 & 0 & 35 & Q4 no baseline; Q6 weak; Q8 none; Q9 child health data, ethics not discussed; Q10 no repetitions. & Abstract and Results (pp.13-19): stability of LIME interpretations between imbalanced and SMOTE-balanced nutritional-status data; read: Publisher PDF via journal viewer (Indonesian Journal of Infomatics 1(2), 2026, open access) \\
\textbf{S56} & Authors: construct validity of SHAP, temporal autocorrelation within students, two datasets only (pp.32-33). & Code archived on Zenodo (doi 10.5281/\allowbreak{}zenodo.22715077); public datasets. & 1 & 1 & 1 & 1 & 1 & 1 & 1 & 1 & .5 & 1 & 95 & Q9 ethics statements present (not applicable); Q10 seeds and folds reported. & RQ2 and Sec.3 (pp.2,9-16): TreeSHAP ranking stability across folds/\allowbreak{}seeds on tabular engineered features; read: Preprint arXiv:2609.28502 (not peer reviewed) \\
\textbf{S57} & Authors: small number of evaluated models for external validation (p.18). Ours: stability is one of four dimensions; perturbations change predictions. & Code: github.com/\allowbreak{}NattakittiP/\allowbreak{}7World\_\allowbreak{}ML\_\allowbreak{}Benchmark (p.23). & 1 & 1 & 1 & 1 & 1 & 1 & 1 & 1 & .5 & 1 & 95 & Q9 clinical data use without ethics detail in sections read; Q10 configurations and bootstrap CIs. & Abstract, Secs.I and III.E (pp.1,3-4,8): SHAP explainability stability (E) reported as a separate score across perturbed synthetic clinical worlds (borderline); read: Publisher PDF (IEEE Access 2026, open access) \\
\textbf{S58} & Authors: the stability objective must itself be estimated (abstract). Ours: retraining stability only. & Public benchmarks used unmodified (Data Availability). & 1 & 1 & 1 & 1 & 1 & 1 & 1 & .5 & .5 & 1 & 90 & Q8 data public, code not seen; Q9 HR data context; Q10 bootstrap refits. & Abstract and Secs.3-7 (HTML full text): SHAP global-ranking agreement across bootstrap refits optimised jointly with AUPRC on 5 tabular attrition benchmarks; read: Publisher HTML (Mathematics 14(18):3283, 2026) \\
\textbf{S59} & Authors discuss attack realism and the PGD anomaly (Sec.10-11). Ours: preprint; adversarial (not random) perturbations; LIME not studied. & Code and results archived on Zenodo (doi 10.5281/\allowbreak{}zenodo.20968236) (ref.21). & 1 & 1 & 1 & 1 & 1 & 1 & .5 & 1 & -- & .5 & 89 & Q7 attack realism discussed; Q8 Zenodo; Q9 no personal data; Q10 repeated runs not clearly reported. & Abstract, Secs.3,7-9 (pp.1-20): TreeSHAP stability under adversarial perturbation for RF/\allowbreak{}XGBoost on 4 tabular security datasets; read: Preprint arXiv:2607.01679 (not peer reviewed) \\
\textbf{S60} & Authors: neural classifiers only; tree-based models left for future work (p.6). Ours: perturbations are adversarial and change predictions, so explanation drift is confounded with prediction change. & Public Kaggle datasets; fixed seeds (p.4); no code link found. & 1 & .5 & 1 & .5 & 1 & .5 & .5 & .5 & -- & .5 & 67 & Q2 architecture details brief; Q4 no explainer comparison; Q6 descriptive plots; Q8 data public, no code; Q10 fixed seed, no repetitions. & Abstract, Secs.3.3-4 (pp.1,3-4): SHAP attribution drift under FGSM/\allowbreak{}PGD perturbations on tabular phishing and UNSW-NB15 data; read: arXiv:2602.06395v1 (states accepted at ICAART 2026; doi 10.5220/\allowbreak{}0014360100004052) \\
\textbf{S61} & Authors: public datasets cannot capture organisational context (Sec.8). Ours: number of splits small/\allowbreak{}unclear; anonymised PCA features limit interpretation. & Public datasets; no code link. & 1 & .5 & 1 & .5 & .5 & .5 & 1 & .5 & .5 & .5 & 65 & Q2 limited detail; Q4 no explainer comparison; Q5 single metric; Q6 descriptive; Q9 governance framing; Q10 fixed seed only. & Abstract, Secs.4.4-6 (pp.1,6-9): SHAP stability across train/\allowbreak{}test splits for LR/\allowbreak{}RF/\allowbreak{}GB on two tabular fraud datasets; read: Publisher PDF (Journal of Computer Science and Technology Studies, published 9 Jan 2026) \\
\textbf{S62} & Ours: single dataset; preprint; single explainer. & Code and data archived on Zenodo (doi 10.5281/\allowbreak{}zenodo.21879038) (p.15). & 1 & 1 & 1 & 1 & 1 & 1 & .5 & 1 & -- & 1 & 94 & Q8 Zenodo archive; Q10 prediction-preserving filter and deterministic partitioning. & Abstract, Secs.3-5 (pp.1-22): TreeSHAP local stability under prediction-preserving perturbations (LR/\allowbreak{}DT/\allowbreak{}RF/\allowbreak{}XGBoost) on CICIoT2023; read: Preprint arXiv:2608.10349 (not peer reviewed) \\
\textbf{S63} & Authors: small target-domain dataset (p.6). Ours: regression task outside finance; only five seeds. & Elite data public (WhoScored); NTHU data private; no code link found. & 1 & 1 & 1 & .5 & 1 & 1 & .5 & .5 & 1 & 1 & 85 & Q4 two models, no explainer baseline beyond CIS; Q7 limited; Q8 partial data; Q9 privacy reason for withholding data; Q10 seeds listed. & Abstract, Secs.3.6-4 (pp.1,7-14): SHAP ranking stability across seeds and domains for RF/\allowbreak{}MLP on tabular event data; read: Preprint arXiv:2605.10796 (not peer reviewed) \\
\textbf{S64} & Authors: limitations section discusses scope; future work on other data. Ours: one dataset. & Data openly available (UCI); no code link found. & 1 & 1 & .5 & 1 & 1 & .5 & 1 & .5 & -- & 1 & 83 & Q3 single dataset; Q6 one dataset; Q8 data only; Q10 noise grid and repeated instances. & Abstract, Methods and Results (Europe PMC full text): AKDE-LIME stability and robustness vs LIME, KernelSHAP, TreeSHAP and Anchor across noise levels 5-20\% on five tree models (tabular UCI data); read: Europe PMC full text PMC13014406 (Applied Artificial Intelligence, 2026, open access) \\
\textbf{S65} & Authors discuss weighting choices and a model-smoothness confound (Sec.6). Ours: preprint; prediction preservation not enforced; small datasets. & Public datasets; no code link found. & 1 & 1 & 1 & 1 & 1 & 1 & .5 & .5 & .5 & .5 & 80 & Q7 robustness checks but limited threats discussion; Q8 no code; Q9 AI Act/\allowbreak{}GDPR context only; Q10 noise draws described, seeds not reported. & Abstract, Secs.3-5 (pp.1-20): SHAP/\allowbreak{}LIME explanation stability index under business noise on 3 tabular datasets; read: Preprint arXiv:2603.05024 (not peer reviewed) \\
\textbf{S66} & Ours: single dataset; preprint; perturbation is bootstrap resampling rather than input noise. & No code link found in the full text. & 1 & 1 & 1 & 1 & 1 & 1 & .5 & 0 & -- & 1 & 83 & Q7 limited threats discussion; Q8 no code/\allowbreak{}data release beyond public benchmark; Q10 bootstrap repetitions. & Abstract, Secs.3-6: multicollinearity-induced SHAP/\allowbreak{}LIME instability on UNSW-NB15 across four model families; read: Preprint arXiv:2605.22529 (not peer reviewed) \\
\textbf{S67} & Authors: single retrospective dataset; no external or prospective validation. & Public Kaggle dataset; no code link found. & 1 & 1 & 1 & 1 & .5 & .5 & 1 & .5 & .5 & 1 & 80 & Q5 composite stability score; Q6 single case study; Q9 clinical data, ethics brief; Q10 30 runs. & Abstract and Methods (HTML full text): explanation stability reported as a separate dimension (30 repeated runs) in a five-dimension trustworthiness framework on tabular METABRIC (borderline); read: Publisher HTML (Future Business Journal 12(1):305, 2026, open access) \\
\textbf{S68} & Authors: bootstrap cost; extension to model-level randomness left for future work (p.9). & Code: github.com/\allowbreak{}Lanxin-Xiang/\allowbreak{}RobustSHAP (p.1). & 1 & 1 & .5 & .5 & 1 & 1 & .5 & 1 & -- & 1 & 83 & Q3 dataset details partly in appendix; Q4 compared with standard SHAP only; Q10 bootstrap replicates. & Abstract, Secs.3-5 (pp.1-9): SHAP ranking variability across resampling/\allowbreak{}seeds and a robust metric (RoSHAP); read: Preprint arXiv:2605.15154 (not peer reviewed) \\
\textbf{S69} & Authors: no protected attributes; simplified profit model (p.15). Ours: only 5 bootstrap replicas; stability is secondary to policy analysis. & Described as reproducible; no code link found. & 1 & 1 & 1 & 1 & 1 & .5 & 1 & .5 & 1 & .5 & 85 & Q6 few replicas; Q8 no code link; Q9 fairness analysed; Q10 B=5 only. & Abstract, Sec.3 (pp.1,5-6): bootstrap SHAP stability is one of four evaluated components on tabular FICO HELOC (borderline); read: Publisher PDF (Artificial Intelligence and Machine Learning Review, 2026; peer-review process not stated in PDF) \\
\textbf{S70} & Authors: threats to validity discussed (Sec.V). Ours: binary malware features; no prediction-preserving filter. & Public malware datasets; code availability not checked in sections read. & 1 & 1 & 1 & 1 & 1 & 1 & 1 & .5 & -- & .5 & 89 & Q8 not verified; Q10 repetitions not fully described in sections read. & Abstract and Secs.III-IV (pp.1,4-9): stability, robustness and effectiveness of LIME, SHAP, Anchor, LORE and LEMNA on five Android-malware feature datasets; read: arXiv:2008.05895 version of the IEEE TIFS article (doi 10.1109/\allowbreak{}TIFS.2020.3021924) \\
\textbf{S71} & Authors: 'nearly optimal' defined by loss; extension to broader models planned (p.11). Ours: global importance only; one dataset. & Not reported in sections read. & 1 & 1 & 1 & .5 & .5 & .5 & 1 & 0 & .5 & .5 & 65 & Q4 compared with single model only; Q5 interval-based; Q6 one dataset; Q8 none found; Q9 ethics approval context; Q10 sampling of models described, seeds not. & Abstract and Secs.2-3 (pp.1-8): SHAP importance across nearly optimal models with uncertainty intervals (ShapleyVIC) on tabular EHR data; read: Preprint arXiv:2212.08370 (no peer-reviewed version identified) \\
\textbf{S72} & Authors: other tree models (e.g., random forest) left for future work (p.6). Ours: short conference paper; neural-network black boxes only. & Code: github.com/\allowbreak{}nranjbar/\allowbreak{}interpretable\_\allowbreak{}machine\_\allowbreak{}learning (p.7). & 1 & .5 & .5 & 1 & .5 & .5 & .5 & 1 & -- & .5 & 67 & Q2/\allowbreak{}Q3 brief; Q5 stability score loosely defined; Q6 small experiments; Q10 repetitions not detailed. & Abstract and Secs.III-IV (pp.1-6): stability, local fidelity and interpretability of tree-ALIME vs ALIME on four tabular datasets (borderline: stability is one of three named outcomes); read: arXiv:2204.03321 version of the CSICC 2022 paper (doi 10.1109/\allowbreak{}CSICC55295.2022.9780503) \\
\textbf{S73} & Authors: binary features and product distributions only; LIME/\allowbreak{}KernelSHAP left open (Conclusion). & Code by collaborators acknowledged; repository not checked. & 1 & 1 & 1 & .5 & 1 & .5 & 1 & .5 & -- & .5 & 78 & Q4 no explainer comparison; Q6 small case study; Q8 not verified; Q10 5 repetitions. & Abstract and Secs.3-5 (pp.1-7): SHAP-score intervals under uncertainty of the feature distribution, with a California Housing decision-tree case study; read: arXiv:2401.12731 version of the ECAI 2024 paper (doi 10.3233/\allowbreak{}FAIA240586) \\
\textbf{S74} & Authors: constraint f\_\allowbreak{}d = phi\_\allowbreak{}d not enforced (Sec.5). Ours: self-explaining model rather than post-hoc SHAP. & Not checked in sections read. & 1 & 1 & 1 & 1 & .5 & 1 & 1 & .5 & -- & 1 & 89 & Q5 uncertainty rather than stability; Q8 not verified; Q10 standard errors reported. & Abstract and Secs.3-5 (pp.1-15): probabilistic Shapley inference giving attribution distributions and credible intervals on synthetic and 8 real tabular datasets; read: Preprint arXiv:2402.04211 (not peer reviewed) \\
\textbf{S75} & Authors: within-neural instability not analysed (p.12). Ours: single-author preprint; disagreement across classes rather than perturbation stability. & Reliability-score code shown in appendix; no repository link found. & 1 & 1 & 1 & 1 & 1 & 1 & 1 & .5 & -- & 1 & 94 & Q8 code snippets only; Q10 multiple seeds controlled. & Abstract, Secs.3-5 (pp.1-7): SHAP/\allowbreak{}LIME disagreement between prediction-equivalent models of different hypothesis classes on 24 tabular datasets, with theory; read: Preprint arXiv:2603.15821 (not peer reviewed) \\
\textbf{S76} & Authors: higher expressiveness can reduce regularity (p.13). & Code: github.com/\allowbreak{}aymen20002005/\allowbreak{}lime\_\allowbreak{}ndt (p.12). & 1 & 1 & .5 & 1 & 1 & 1 & .5 & 1 & -- & 1 & 89 & Q3 dataset details in tables; Q7 brief; Q10 repeated perturbation runs. & Abstract and Secs.2.2-4 (pp.1,4-13): stability, fidelity and regularity of NDT-LIME vs LIME surrogates on benchmark tabular datasets (borderline); read: Preprint arXiv:2603.20919 (not peer reviewed) \\
\textbf{S77} & Ours: simulations only; single-author preprint; argues against seeking more stable explainers. & No code link found. & 1 & 1 & .5 & 1 & 1 & .5 & 1 & 0 & -- & 1 & 78 & Q3 synthetic only; Q6 no real data; Q8 none; Q10 simulation design specified. & Abstract and Secs.3-5 (pp.1-28): explanatory instability of local-linear and Shapley explanations tied to forecast uncertainty near segment/\allowbreak{}decision boundaries, tested in simulations; read: Preprint arXiv:2602.00179v3 (not peer reviewed) \\
\textbf{S78} & Ours: data not shareable; stability protocol briefly described; no repetitions reported. & MIMIC-IV not redistributable; fixed seeds; no code link (p.6). & 1 & .5 & 1 & 1 & 1 & .5 & .5 & 0 & .5 & .5 & 65 & Q2 brief; Q6 single dataset; Q8 none; Q9 data-use agreement; Q10 fixed seeds only. & Abstract and Secs.5-6 (pp.1-10): stability (Lipschitz estimate) of TreeSHAP, ensemble LIME and Anchors on MIMIC-IV mortality models (borderline: one of three criteria); read: Publisher PDF (Int. J. Computer Sciences and Engineering 14(5), 2026, open access) \\
\textbf{S79} & Authors: single event inventory; no spatial-block validation (p.18). & Implementation details in appendix; code not checked. & 1 & 1 & 1 & 1 & 1 & 1 & 1 & .5 & -- & 1 & 94 & Q8 not verified; Q10 repeated seeds. & Abstract and Secs.II-IV (pp.1-8): SHAP stability under Gaussian perturbation across DEM resolutions for LR, SVM, RF and XGBoost with repeated seeds (tabular slope-unit factors); read: Publisher PDF (IEEE JSTARS 19, 2026, open access) \\
\end{longtable}
